\documentclass[../../main2.tex]{subfiles}

\begin{document}

	% 全书位置：导论 → 动机 → 【目标】 → 预测与因果 → ……
	% v2 要点：本章零公式（打分表代替公式盒）；新增 §2.4「看得见的约束」
	%（原 v1 Part II 状态空间层的降级归宿）；章末埋「微观→宏观」之问（Part IV 开头收）。

	\textbf{\large 全书位置}：动机 → \textbf{目标} → 预测与因果。

	\textit{把「想要」从一个名词改造成一个排序；给排序打分加权；再诚实地承认权重看不见——最后清点预测者真正的原料。}

	\vspace{1em}

	\chapter{目标：把意向写成可比较的量}
	\label{ch:goal}

	\begin{motivation}
		上一章遗留的问题：动机是预测的必要输入，但从动机到结果隔着「目标」与「环境」两道转换。环境下一节就有安排；本章处理前一层——但「目标」这个词还只是日常词汇：「我想要 X」，X 到底是什么形状的东西？一个名词？一张清单？还是一个数？\\
		\textbf{核心动机}：把「目标」从愿望改写成可比较、可打分、可从行为反推的量。没有这一步，「动机与结果的落差」就只是修辞，预测就无从谈起。
	\end{motivation}

	\begin{logicmap}
	\begin{tikzpicture}[node distance=0.9cm, every node/.style={draw, rectangle, rounded corners, align=center}]
	\node (q) {「想要 X」：X 是什么？};
	\node (rank) [below=of q] {目标 = 排序（§2.1）};
	\node (comm) [below=of rank] {尺子不止一把：不可通约（§2.2）};
	\node (util) [below=of comm] {给想要的东西打分（§2.3）};
	\node (con) [below=of util] {看得见的约束：菜单与点菜（§2.4）};
	\node (rev) [below=of con] {权重不可见，从行为反推（§2.5）};
	\node (next) [below=of rev] {原料齐了，规律会骗人（§2.6）};
	\draw[-{Stealth}] (q) -- (rank);
	\draw[-{Stealth}] (rank) -- (comm);
	\draw[-{Stealth}] (comm) -- (util);
	\draw[-{Stealth}] (util) -- (con);
	\draw[-{Stealth}] (con) -- (rev);
	\draw[-{Stealth}] (rev) -- (next);
	\end{tikzpicture}
	\end{logicmap}

	% ==================================================================
	\section{目标不是名词，是排序}
	\label{sec:goal-ranking}

	\begin{readingnote}
		你有没有想过，为什么「我的目标是有钱、有闲、有健康」这种话，说的时候很爽，说完等于没说？\\
		因为这三样东西\textbf{打架}的时候你听谁的，才是你的目标——清单本身不是。
	\end{readingnote}

	先做一个残酷但诚实的观察：没有人只「想赚钱」。一个人或者想要「钱多过现在的两倍」，或者想要「钱够花但六点下班」，或者想要「钱比同学多」。「想赚钱」不是一个目标，是一个\textbf{目标集合}——而集合不驱动行为，取舍才驱动行为。

	场景永远是取舍：加班（钱更多、健康更糟、闲暇更少）还是回家？跳槽（成长加分、稳定减分）还是留下？一个人在冲突时刻的\textbf{顺序}，才是他的目标。所以：

	\begin{quote}
		目标的正确形状不是名词，是\textbf{排序（preference ordering）}——在任意两个选项之间，排出先后。
	\end{quote}

	用本书的标准句式把这个概念分离出来：

	\begin{quote}
		如果「想要」不需要比较，那它只是情绪——一阵想要，一阵不想要，预测不了任何行为。但现在加入一个约束：\textbf{目标必须能指导选择}——能指导选择的东西，必须能在任意两个选项之间排出先后。这一步改变了所有性质，这就是「目标」从「愿望」中分离出来的时刻。
	\end{quote}

	\textit{生活类比}：买房的人嘴上说「要近地铁、要好学区、要大平米、要便宜」——全都要。但预算一压，中介只带他看两套房：A 近地铁 60 平，B 远地铁 100 平学区好。他选了哪个，他的「目标」就写在哪个选择里。嘴上的清单是广告，行为上的排序才是本人。

	要当「排序」，得满足两条朴素的品质。第一，\textbf{没有「说不清」}：任何两个选项摆在一起，总能给出先后（允许打平）——因为现实迟早会把这两选项端到你面前逼你选。第二，\textbf{没有「循环」}：不能面比饭好、饭比面包好、转一圈又面包比面好——有循环的人，谁先骗他一票，他就要什么。满足这两条，一个人才谈得上「有目标」；也只有满足这两条，别人才可能从他的行为里读出稳定的偏好来——这件事 §\ref{sec:goal-reveal} 要用。

	\begin{questionbox}
		\textit{现在你可能想反驳：「排序太苛刻了吧？我很多时候真的无所谓——中午吃面还是吃饭，随便是真心话。」}\\
		\textbf{无所谓本身就是排序的一部分}：打平是合法的平局，不破坏第一条品质。真正破坏排序的是\textbf{循环偏好}：面排在饭前面，饭排在面包前面，绕一圈面包又排回面前面。行为经济学证明循环偏好在人类身上真实存在（框架效应（framing effect）可以把偏好当场反转）。\lowfreq{本书对这一事实的处理不是假装它不存在，而是接受一个较弱的工作假设：在生存相关的重要取舍上，偏好近似可排序；不成立的情况，本书会如实列出，而不假装它们不存在。}
	\end{questionbox}

	\paragraph*{逻辑衔接}
	排序解决了「能不能比」。但「比」需要尺子——所有选项都放在同一把尺子上量吗？下一节给出坏消息：尺子不止一把，而且有些尺子之间\textbf{没有换算率}。

	% ==================================================================
	\section{可通约性难题}
	\label{sec:goal-commensurability}

	\begin{readingnote}
		你有没有想过，「这条命值多少钱」这个问题，为什么让人本能地不舒服？\\
		不是问题残忍——是你心里知道：\textbf{有些东西换算不成别的}。这不是道德洁癖，是目标世界的结构。
	\end{readingnote}

	三样东西放在一起：金钱、自由、生命。经济学教科书喜欢假装它们都能换算——一切皆有价。但做一个思想实验就能摸到边界：保险公司确实为生命定价（统计生命价值，value of a statistical life），但注意它定的是什么——是\textbf{降低 $0.1\%$ 死亡风险}的价格，不是\textbf{生命}的价格。你愿意花两千块买保险降低 $0.1\%$ 的风险，和你愿意收多少钱\textbf{去死}，是两个宇宙里的问题。一旦越过某条线，「报价」这个动作本身失灵。

	现在给出本书在这一节的明确回答——不是折中，是立场：

	\begin{quote}
		\textbf{目标之间不能完全通约。且这不是人类理性的缺陷——这恰恰是政治与道德存在的理由。}
	\end{quote}

	反过来想就通了：如果生命、自由、金钱之间真的存在完备的换算表，那么一切价值冲突都只是计算错误，人类就不需要道德争论，只需要\textbf{会计}。正因为生命不能完全换成自由、自由不能完全换成钱，排序冲突才无解于算盘，人类才需要一样超出计算的装置来裁决——伦理是它，政治也是它（第\ref{ch:politics}章将看到政治如何接手这个裁决）。

	\begin{quote}
		如果所有目标都能换算成单一尺度，目标世界就退化成一条线——世界只需要会计学。但现在加入一个约束：\textbf{某些目标之间的兑换率不存在，或视之为无限大}。这一步改变了所有性质：目标从一条线变成一张多维的清单，排序只能逐项进行、局部进行，全局排序必须引入权重——这就是多目标问题的诞生时刻。
	\end{quote}

	「不能完全换算」在书面上叫\textbf{不可通约（incommensurable）}——此后本书只用日常说法：\textbf{尺子不止一把}。它的直接后果是：有些选项之间真的比不出全局胜负——A 这项好、B 那项好，谁也不压倒谁。这个「比不出的地带」，就是价值冲突的地盘，也是下一节「打分加权」要出场的原因：现实逼我们行动，行动逼我们把多维压回一维。

	\begin{questionbox}
		\textit{现在你可能想反驳：「经济学不就假设人人效用最大化吗？那不就是假定一切都通约了？」}\\
		\textbf{把两件事分开：追求「最好」和「一切向钱看」是两回事。}「理性人」假设的是\textbf{排序的一致性}（上一节的两条品质），不是「万物皆可换钱」。把经济学读成拜金教程，是把「打分加权」误认成了「统一币制」。
	\end{questionbox}

	\begin{reader_anticipation}
		\item 不可通约的排序冲突，最后由谁裁决？→ 第\ref{ch:politics}章「权力」：政治正是交换中加入不可定价的偏好之后长出的东西。
		\item 多个因素搅在一起「各占多少」时，会不会也遇到「加不起来」的麻烦？→ 会。第\ref{ch:contribution}章与第\ref{ch:law-method}章：推手之间的交互，就是不可通约在因果世界的影子。
	\end{reader_anticipation}

	\paragraph*{逻辑衔接}
	尺子不止一把，可现实中的选择却总是「二选一」——多维的想要必须被压回一维才能行动。这个压缩装置就是下一节的主角：打分加权。它是工程奇迹，也是本章所有麻烦的起点。

	% ==================================================================
	\section{目标函数：给想要的东西打分}
	\label{sec:goal-utility}

	\begin{readingnote}
		你有没有想过，为什么招股书、征婚广告、择校清单，最后都变成一张「打分加权表」？\\
		因为多维的想要，人类大脑算不动——只能压成一维再选。这个压缩动作里，藏着\textbf{权重}的秘密。
	\end{readingnote}

	工程答案先摆出来：给每把尺子配一个权重，把多维压回一个总分。这就是经济学里的\textbf{目标函数（utility function，效用函数）}——先用日常版：一张打分表。

	一个具体的例子：选工作。三份 offer，四个维度，求职者心里那张表大概长这样：

	\begin{table}[h]
	\centering
	\begin{tabular}{lcccc}
		\toprule
		 & 钱 & 通勤 & 成长 & 稳定 \\
		\midrule
		甲（大厂） & 9 & 3 & 5 & 8 \\
		乙（创业） & 5 & 7 & 9 & 2 \\
		丙（国企） & 6 & 8 & 4 & 9 \\
		\bottomrule
	\end{tabular}
	\caption{打分表：每列一把尺子，各行是选项。}
	\end{table}

	表里还没有权重。现在把权重请进来——钱占几成、成长占几成——每行加权求和，得到三个总分，最高的那个就是「他的选择」。他最后说「我选乙」的那一刻，其实是报出了一个加权结果：成长在他心里的权重，压过了钱和稳定。

	\textit{生活类比}：这张表像导航软件。城市有千万条路、无数个红绿灯和限速杆，导航把它压缩成一张几寸屏上的图——丢掉了九成的真实，保住了你要做决定的那一成。$1:1$ 的导航是没用的：它得装下整座城市。打分表对「想要」干的就是这件事：压缩必然丢信息，\textbf{丢什么，由权重决定}。\footnote{博尔赫斯在短篇《论科学的精确性》里写过一张 $1:1$ 的帝国地图——铺满整个帝国，被风沙埋掉。它常被拿来讽刺「不压缩的认知」，此处留个脚印。}

	现在是本章最重要的一段诚实话：\textbf{权重不可观测}。没有人脑门上写着「健康 $0.3$、钱 $0.4$」。更糟的是权重不稳定——同一个人，体检报告出来那天，健康的权重飙升；季度奖金公布那天，钱的权重飙升。所以打分表是「结构上存在、数值上悬空」的量表。\lowfreq{目标函数是建模约定，不是定理；权重是自由参数，不是测量值。}

	\begin{questionbox}
		\textit{现在你可能想反驳：「既然权重不可观测，这套打分不就是装饰品吗？」}\\
		\textbf{部分正确——所以要看清它到底贡献了什么。}打分的作用不是被填满，而是保证「选择有结构」：只要那张表在且不瞬变，行为就不是纯随机，就\textbf{可被统计检验}。真正可算的东西不在表里，而在表外——打分表的影子会从行为里漏出来。下一节就去找这个影子。
	\end{questionbox}

	\paragraph*{逻辑衔接}
	打分的语法有了，权重却悬在半空。但先别急着处理它——还有一类东西比权重更实在、也更容易被忽略：那些\textbf{看得见的约束}。它们决定选项清单本身。下一节给它们正式的位置。

	% ==================================================================
	\section{看得见的约束：菜单与点菜}
	\label{sec:goal-constraints}

	\begin{readingnote}
		你有没有想过，为什么沙盘上的河流与山脉，能决定两千年后谁修运河、谁建舰队？\\
		地理不说话——但看起来，它早就把答案写在了地形上。这是真的吗？
	\end{readingnote}

	先把话说小：预测者眼里的世界，清单其实很短——\textbf{看得见的约束}（地理、资源、人口、技术存量），\textbf{藏着动机的人}，以及这两者\textbf{没完没了的推挤}。本节处理第一样：看得见的约束。它是三章之前被否掉的那层「本体论」里，唯一真正值得留下的部分——不是作为哲学前提，而是作为\textbf{工具}。

	看得见的约束做什么用？一句话：\textbf{它决定菜单，不决定点菜}。

	港口城市长得像港口城市：码头、船坞、外贸行、移民社区；绿洲城市长得像绿洲城市：水权法庭、骆驼市集、驿站。黄河流域给它的居民上的菜单上写着「治水协作」，爱琴海的菜单上写着「跨海贸易」。上一章那对分岔——大一统与城邦——不是民族的性格差异，是两份不同的菜单上，各自被点了的菜。

	点菜的是谁？是拿着各自打分表的人。菜单管住的是\textbf{可行}：沙漠里点不出水稻文明，岛上点不出陆权帝国；点菜管住的是\textbf{选择}：同一份黄河菜单上，也有人点盐铁专营、有人点黄老无为。所以本书把地理决定论降级安置在这里：\textbf{约束是预测的第一味原料，不是原因的终点}——把菜单当成食客，就又回到了「数据最多、预测最离谱」的老路。

	菜单上还有一条特殊的底线，值得单独说：所有打分表共享同一个定义域——\textbf{存在}。赌桌上没有「下一局」的人，谈不上翻盘；把「我们自己不存在」放进选项清单，一切排序立刻失效。所以任何目标函数都默认一条不可谈判的约束：\textbf{别把桌子掀了}。这不是道德律令，是排序能成立的前提。

	最后一句提醒：菜单本身也在变。技术会重画地图——运河会淤，铁路会绕开港口，互联网把「距离」这个千年约束改小了一号。看得见的约束是慢变量，不是常量。

	\begin{questionbox}
		\textit{现在你可能想反驳：「决定菜单不决定点菜？那不就是半吊子的地理决定论。」}\\
		\textbf{差别在于预测时你往模型里装什么。}纯决定论装地形，然后硬推；菜单思维装「地形给定的选项清单+人的打分表」，然后问：这份菜单上，最可能被点的是哪道菜？后者多装的那一味，就是第1章说的必要输入——动机。
	\end{questionbox}

	\begin{reader_anticipation}
		\item 菜单（约束）怎么进入后面的工具箱？→ 它是「无你检验」里那个被按住的背景——第\ref{ch:contribution}章。
		\item 技术重画菜单，有没有正式的处理？→ 第\ref{ch:forecast}章的失败模式清单里，「外部冲击」一条专治这个。
	\end{reader_anticipation}

	\paragraph*{逻辑衔接}
	看得见的约束画好了边界，菜单摆上了桌。但桌上真正点菜的东西——那张藏在她/他心里的打分表——看不见。看不见的东西怎么研究？下一节：从「做了什么」反推「想要什么」。

	% ==================================================================
	\section{目标函数如何反推：显示性偏好}
	\label{sec:goal-reveal}

	\begin{readingnote}
		你有没有想过，为什么「听其言」之后还要「观其行」？\\
		因为嘴上的目标可以无限美好，行为却会把真实的排序\textbf{供}出来。
	\end{readingnote}

	经济学家萨缪尔森把这个直觉锻造成了工具：\textbf{显示性偏好（revealed preference）}。逻辑只有一句话：一个人在能选甲的情况下选了乙，说明\textbf{对他而言}乙不比甲差。不需要读心，只需要记账——账本就是他的选择史。

	例子：同事口头宣称「健康最重要」，但连续三年接受每天通勤三小时的高薪岗位。他的嘴和脚给出两份不同的权重表。显示性偏好只认脚。

	\begin{quote}
		如果权重可观测，反推就无必要——直接读数即可。但现在加入一个约束：\textbf{权重藏在心里}。这一步改变了所有性质：目标从「给定」变成「推断」，从描述问题变成统计问题——从此，「想要什么」永远隔着数据的一层纱。
	\end{quote}

	然后是本章第二个重要的诚实话：\textbf{反推有唯一性问题}。同一组行为，可以被多张不同的打分表解释。买咖啡这个行为，与三种互不相同的人都兼容：真喜欢咖啡味道的、需要咖啡因提神的、办公室文化里不想显得不合群的。三种动机，同一张消费小票。行为数据无法区分它们——这不是数据不够多，是\textbf{结构性的分不清}。

	单个行为只能给出偏好序的\textbf{下界}：它排除了「被否决的选项排在所选选项之上」，仅此而已。要收窄，只有一条路：让同一个人在不同约束下反复选择——价格变、环境变、同伴变——每次选择切掉一块候选打分表的 space。但无论切多少刀，有限次观测永远定不出唯一的那张表。

	\begin{questionbox}
		\textit{现在你可能想反驳：「行为经济学早就证明偏好会被框架效应反转——反推出来的排序还作数吗？」}\\
		\textbf{作数，但要换个读法：局部偏好。}反推出的排序只在\textbf{观测到的约束区域内}有效，出了这个区域不作数。偏好是有背景的——这话不是让步，是预告：有背景的偏好，恰好是第\ref{ch:prediction}章「相关与因果」问题的常客。
	\end{questionbox}

	\begin{reader_anticipation}
		\item 反推不唯一，「各占多少」还怎么算？→ 第\ref{ch:contribution}章：无你检验不要求唯一识别全部动机，只要求对比「有 A / 无 A」两个世界——天然绕开唯一性。
	\end{reader_anticipation}

	\paragraph*{逻辑衔接}
	现在清点本章家当：目标是排序（§\ref{sec:goal-ranking}）；尺子不止一把（§\ref{sec:goal-commensurability}）；打分加权但权重悬空（§\ref{sec:goal-utility}）；看得见的约束画菜单（§\ref{sec:goal-constraints}）；从行为能反推、但不唯一（§\ref{sec:goal-reveal}）。把这些加起来，预测者的原料清单就齐了。下一节收束本章——并把一个更大的问题按下不表。

	% ==================================================================
	\section{本章小结：原料齐了，但规律会骗人}
	\label{sec:goal-summary}

	\begin{readingnote}
		你有没有想过，为什么菜谱齐全、食材新鲜，做出来的菜还是可能难吃？\\
		因为\textbf{原料齐}和\textbf{会做}之间，还差一门手艺。预测的手艺，下一部分才开始教。
	\end{readingnote}

	Part I 到此收口。预测者手里现在有什么？
	\begin{itemize}
		\item \textbf{动机}：预测的必要输入，一切模型的起点（第1章）；
		\item \textbf{目标}：动机的可比较形状——排序、多把尺子、一张悬空的打分表（本章）；
		\item \textbf{看得见的约束}：菜单——可行性的硬边界（本章 §\ref{sec:goal-constraints}）；
		\item \textbf{行为}：动机留下的唯一脚印，可观测、可反推、但不唯一（本章 §\ref{sec:goal-reveal}）。
	\end{itemize}

	接下来的一步，看起来朴素得不像话：拿着这些原料，\textbf{从过去的行为里找规律}——谁跟谁一起涨、什么之后总跟着什么——然后把规律外推到明天。这几乎就是常识本身。

	但导论里那群最会算的人，手里拿的正是这份原料，做的正是这个动作，然后集体撞墙。问题不在原料，在\textbf{规律}这个词的歧义：规律有两种，一种真的能带你到明天，一种只会把你留在昨天——而它们在数据表上\textbf{长得一模一样}。怎么把这两种规律区分开，是第\ref{ch:prediction}章的全部内容。

	最后，按约定埋下一个问题，先按下不表：\textbf{一个人的动机好懂，一百万人的动机会汇成什么？}一个人想赚钱，不等于市场有规律；一个人怕死，不等于社会怕死。个人的「想要」和宏观的「大势」之间隔着什么——这个问题会在第\ref{ch:channels}章（四大齿轮的入口）被正面回答。现在只需要记住：我们迟早要过这座桥。

	\paragraph*{逻辑衔接}
	原料清点完毕，陷阱已预告。第\ref{ch:prediction}章「预测与因果」将从一次人类历史上最贵的「规律失灵」讲起——看看拥有数据最多的人，是怎么被自己的规律骗的。

\end{document}
